\documentclass[12pt,a4paper]{article}
\usepackage{geometry}
\geometry{left=2.5cm,right=2.5cm,top=2.0cm,bottom=2.5cm}
\usepackage[english]{babel}
\usepackage{amsmath,amsthm}
\usepackage{amsfonts}
\usepackage[longend,ruled,linesnumbered]{algorithm2e}
\usepackage{fancyhdr}
\usepackage{ctex}
\usepackage{array}
\usepackage{listings}
\usepackage{color}
\usepackage{graphicx}

\begin{document}


\title{
{《优化理论与算法》第 {$4$} 次作业
}
}
\date{}

 \author{
 姓名：\underline{焦传斌}~~~~~~
 学号：\underline{2024111223}~~~~~~}

\maketitle

\begin{itemize}
	\item 作业截止于{\bf 周五（4/3/2026）}，在Canvas系统中提交PDF或照片文件
	\item 作业题目的tex文件可以在Canvas中下载
	\item 鼓励相互讨论，但不要抄袭.
\end{itemize}
%%%%%%%%%%%%%%%%%%%%%%%%%%%%%%%%%%%%%%%%%%%%%%%%%%%%%%%%%%%%%%%%%
%%%%%%%%%%%%%%%%%%%%%%%%%%%%%%%%%%%%%%%%%%%%%%%%%%%%%%%%%%%%%%%%%

\section{阅读与积累}
\paragraph{1.1}
\begin{itemize}
	\item 阅读Bertsimas and Tsitsikils，也就是Introduction to Linear Optimization教材，第3章与第4.1-4.4节相关内容。
\end{itemize}

\section{习题}


\paragraph{2.1}
在利用表格法求解线性规划过程中，假设当前基为 $B=\{3,4,5\}$有如下表格：
\begin{center}
	\begin{tabular}{|c| c c c c c |c|} 
		\hline
		B & $\delta$ &-2 &0 &0 &0 & -10 \\
		\hline
		3 & -1 & $\eta$ &1 &0 &0 & 4 \\
		4 & $\alpha$ & -4 &0 &1 &0 & 1 \\
		5 & $\gamma$ & 3 &0 &0 &1 & $\beta$ \\
		\hline
	\end{tabular}
\end{center}
其中，$\alpha$, $\beta$, $\gamma$, $\delta$, $\eta$为表中未知参数。请给出下面命题成立时$(\alpha, \beta, \gamma, \delta, \eta)$的取值范围
\begin{enumerate}
	\item 当前解是可行解，但不是最优解
	\item 问题是无界问题
	\item 当前解是最优解
\end{enumerate}

\textbf{解：}

根据单纯形法表格的性质，基变量为$\{3,4,5\}$，非基变量对应第1、2列。

\textbf{(1) 当前解是可行解，但不是最优解：}

可行解要求：右端项非负，即  $\beta \geq 0$

不是最优解：存在负检验数。第2列检验数$-2 < 0$已满足

即：$\beta \geq 0$，$\alpha, \gamma, \delta, \eta$ 任意

\textbf{(2) 问题是无界问题：}

无界条件：存在负检验数的非基变量列，其所有系数都$\leq 0$

第2列检验数$-2 < 0$，该列系数为$(\eta, -4, 3)^T$，但$3 > 0$

对于第1列，若$\delta < 0$，该列系数为$(-1, \alpha, \gamma)^T$

无界条件：$\delta < 0$, $-1 \leq 0$, $\alpha \leq 0$, $\gamma \leq 0$

同时需保证可行：$\beta \geq 0$

取值范围：$\delta < 0, \alpha \leq 0, \gamma \leq 0, \beta \geq 0, \eta \text{ 任意}$

\textbf{(3) 当前解是最优解：}

最优条件：可行且所有检验数$\geq 0$

可行性：$\beta \geq 0$

最优性：第2列检验数$-2 < 0$，不满足

所以无解。

\paragraph{2.2}
利用两阶段法求解如下线性规划

\begin{equation*}
\begin{array}{lllll@{}ll}
\text{minimize}  & -3x_1&+x_2&-2x_3 &\\
\text{subject to}& x_1 &+3x_2  &+x_3  &\leq 5 \\
&2x_1 &-x_2 & +x_3  & \geq 2  \\
&4x_1 &+3x_2 &-2x_3   & =5  \\
&x_1 \geq 0, &x_2 \geq 0, & x_3 \geq 0 &  
\end{array}
\end{equation*}

\textbf{解：}

\textbf{标准形式：}引入松弛变量$x_4$，剩余变量$x_5$，人工变量$x_6, x_7, x_8$
\begin{align*}
& x_1 + 3x_2 + x_3 + x_4 + x_6 = 5 \\
& 2x_1 - x_2 + x_3 - x_5 + x_7 = 2 \\
& 4x_1 + 3x_2 - 2x_3 + x_8 = 5
\end{align*}

\textbf{第一阶段：}$\min w = x_6 + x_7 + x_8$

初始单纯形表（基$B = \{x_6, x_7, x_8\}$）：

\begin{center}
\begin{tabular}{|c|rrrrrrrr|r|}
\hline
B & $x_1$ & $x_2$ & $x_3$ & $x_4$ & $x_5$ & $x_6$ & $x_7$ & $x_8$ & $b$ \\
\hline
 & -7 & -5 & 0 & -1 & 1 & 0 & 0 & 0 & -12 \\
\hline
$x_6$ & 1 & 3 & 1 & 1 & 0 & 1 & 0 & 0 & 5 \\
$x_7$ & 2 & -1 & 1 & 0 & -1 & 0 & 1 & 0 & 2 \\
$x_8$ & 4 & 3 & -2 & 0 & 0 & 0 & 0 & 1 & 5 \\
\hline
\end{tabular}
\end{center}

$x_1$入基，比值检验：$\min\{5/1, 2/2, 5/4\} = 1$，$x_7$出基。

第1次迭代后（基$B = \{x_6, x_1, x_8\}$）：

\begin{center}
\begin{tabular}{|c|rrrrrrrr|r|}
\hline
B & $x_1$ & $x_2$ & $x_3$ & $x_4$ & $x_5$ & $x_6$ & $x_7$ & $x_8$ & $b$ \\
\hline
 & 0 & $-\frac{17}{2}$ & $\frac{7}{2}$ & -1 & $-\frac{5}{2}$ & 0 & $\frac{7}{2}$ & 0 & -5 \\
\hline
$x_6$ & 0 & $\frac{7}{2}$ & $\frac{1}{2}$ & 1 & $\frac{1}{2}$ & 1 & $-\frac{1}{2}$ & 0 & 4 \\
$x_1$ & 1 & $-\frac{1}{2}$ & $\frac{1}{2}$ & 0 & $-\frac{1}{2}$ & 0 & $\frac{1}{2}$ & 0 & 1 \\
$x_8$ & 0 & 5 & -4 & 0 & 2 & 0 & -2 & 1 & 1 \\
\hline
\end{tabular}
\end{center}

$x_2$入基，比值检验：$\min\{4/(7/2), 1/5\} = 1/5$，$x_8$出基。

第2次迭代后（基$B = \{x_6, x_1, x_2\}$）：

\begin{center}
\begin{tabular}{|c|rrrrrrrr|r|}
\hline
B & $x_1$ & $x_2$ & $x_3$ & $x_4$ & $x_5$ & $x_6$ & $x_7$ & $x_8$ & $b$ \\
\hline
 & 0 & 0 & $-\frac{33}{10}$ & -1 & $\frac{9}{10}$ & 0 & $\frac{1}{10}$ & $\frac{17}{10}$ & $-\frac{33}{10}$ \\
\hline
$x_6$ & 0 & 0 & $\frac{33}{10}$ & 1 & $-\frac{9}{10}$ & 1 & $\frac{9}{10}$ & $-\frac{7}{10}$ & $\frac{33}{10}$ \\
$x_1$ & 1 & 0 & $\frac{1}{10}$ & 0 & $-\frac{3}{10}$ & 0 & $\frac{3}{10}$ & $\frac{1}{10}$ & $\frac{11}{10}$ \\
$x_2$ & 0 & 1 & $-\frac{4}{5}$ & 0 & $\frac{2}{5}$ & 0 & $-\frac{2}{5}$ & $\frac{1}{5}$ & $\frac{1}{5}$ \\
\hline
\end{tabular}
\end{center}

$x_3$入基，比值检验：$\min\{(33/10)/(33/10), (11/10)/(1/10)\} = 1$，$x_6$出基。

第3次迭代后（基$B = \{x_3, x_1, x_2\}$），第一阶段达到最优：

\begin{center}
\begin{tabular}{|c|rrrrrrrr|r|}
\hline
B & $x_1$ & $x_2$ & $x_3$ & $x_4$ & $x_5$ & $x_6$ & $x_7$ & $x_8$ & $b$ \\
\hline
 & 0 & 0 & 0 & 0 & 0 & 1 & 1 & 1 & 0 \\
\hline
$x_3$ & 0 & 0 & 1 & $\frac{10}{33}$ & $-\frac{3}{11}$ & $\frac{10}{33}$ & $\frac{3}{11}$ & $-\frac{7}{33}$ & 1 \\
$x_1$ & 1 & 0 & 0 & $-\frac{1}{33}$ & $-\frac{3}{11}$ & $-\frac{1}{33}$ & $\frac{3}{11}$ & $\frac{4}{33}$ & 1 \\
$x_2$ & 0 & 1 & 0 & $\frac{8}{33}$ & $\frac{2}{11}$ & $\frac{8}{33}$ & $-\frac{2}{11}$ & $\frac{1}{33}$ & 1 \\
\hline
\end{tabular}
\end{center}

第一阶段$w^* = 0$，原问题可行。删去人工变量列，进入第二阶段。

\textbf{第二阶段：}$\min z = -3x_1 + x_2 - 2x_3$

初始单纯形表（基$B = \{x_3, x_1, x_2\}$）：

\begin{center}
\begin{tabular}{|c|rrrrr|r|}
\hline
B & $x_1$ & $x_2$ & $x_3$ & $x_4$ & $x_5$ & $b$ \\
\hline
 & 0 & 0 & 0 & $\frac{3}{11}$ & $-\frac{17}{11}$ & -4 \\
\hline
$x_3$ & 0 & 0 & 1 & $\frac{10}{33}$ & $-\frac{3}{11}$ & 1 \\
$x_1$ & 1 & 0 & 0 & $-\frac{1}{33}$ & $-\frac{3}{11}$ & 1 \\
$x_2$ & 0 & 1 & 0 & $\frac{8}{33}$ & $\frac{2}{11}$ & 1 \\
\hline
\end{tabular}
\end{center}

$x_5$入基，比值检验：$\min\{1/(2/11)\} = 11/2$，$x_2$出基。

第1次迭代后（基$B = \{x_3, x_1, x_5\}$），达到最优：

\begin{center}
\begin{tabular}{|c|rrrrr|r|}
\hline
B & $x_1$ & $x_2$ & $x_3$ & $x_4$ & $x_5$ & $b$ \\
\hline
 & 0 & $\frac{17}{2}$ & 0 & $\frac{7}{3}$ & 0 & $\frac{25}{2}$ \\
\hline
$x_3$ & 0 & $\frac{3}{2}$ & 1 & $\frac{2}{3}$ & 0 & $\frac{5}{2}$ \\
$x_1$ & 1 & $\frac{3}{2}$ & 0 & $\frac{1}{3}$ & 0 & $\frac{5}{2}$ \\
$x_5$ & 0 & $\frac{11}{2}$ & 0 & $\frac{4}{3}$ & 1 & $\frac{11}{2}$ \\
\hline
\end{tabular}
\end{center}

所有检验数$\geq 0$，达到最优。

\textbf{最优解：}$x_1 = \dfrac{5}{2}, x_2 = 0, x_3 = \dfrac{5}{2}$

\textbf{最优值：}$z_{\min} = -\dfrac{25}{2}$

\paragraph{2.3}
利用两阶段法求解如下线性规划

\begin{equation*}
\begin{array}{lllll@{}ll}
\text{maximize}  & x_1&-x_2&+x_3 &\\
\text{subject to}& 2x_1 &-x_2  &-2x_3  &\leq 4 \\
&2x_1 &-3x_2 & -x_3  & \leq -5  \\
&-x_1 &+x_2 &+x_3   & \leq -1  \\
&x_1 \geq 0, &x_2 \geq 0, & x_3 \geq 0 &  
\end{array}
\end{equation*}

\textbf{解：}

转化为最小化：$\min z = -x_1 + x_2 - x_3$

\textbf{标准形式：}第2、3个约束右端为负，乘以$-1$后引入松弛变量$x_4$，剩余变量$x_5, x_6$，人工变量$x_7, x_8, x_9$：
\begin{align*}
& 2x_1 - x_2 - 2x_3 + x_4 + x_7 = 4 \\
& -2x_1 + 3x_2 + x_3 - x_5 + x_8 = 5 \\
& x_1 - x_2 - x_3 - x_6 + x_9 = 1
\end{align*}

\textbf{第一阶段：}$\min w = x_7 + x_8 + x_9$

初始单纯形表（基$B = \{x_7, x_8, x_9\}$）：

\begin{center}
\begin{tabular}{|c|rrrrrrrrr|r|}
\hline
B & $x_1$ & $x_2$ & $x_3$ & $x_4$ & $x_5$ & $x_6$ & $x_7$ & $x_8$ & $x_9$ & $b$ \\
\hline
 & -1 & -1 & 2 & -1 & 1 & 1 & 0 & 0 & 0 & -10 \\
\hline
$x_7$ & 2 & -1 & -2 & 1 & 0 & 0 & 1 & 0 & 0 & 4 \\
$x_8$ & -2 & 3 & 1 & 0 & -1 & 0 & 0 & 1 & 0 & 5 \\
$x_9$ & 1 & -1 & -1 & 0 & 0 & -1 & 0 & 0 & 1 & 1 \\
\hline
\end{tabular}
\end{center}

$x_1$入基，比值检验：$\min\{4/2, 1/1\} = 1$，$x_9$出基。

第1次迭代后（基$B = \{x_7, x_8, x_1\}$）：

\begin{center}
\begin{tabular}{|c|rrrrrrrrr|r|}
\hline
B & $x_1$ & $x_2$ & $x_3$ & $x_4$ & $x_5$ & $x_6$ & $x_7$ & $x_8$ & $x_9$ & $b$ \\
\hline
 & 0 & -2 & 1 & -1 & 1 & 0 & 0 & 0 & 1 & -9 \\
\hline
$x_7$ & 0 & 1 & 0 & 1 & 0 & 2 & 1 & 0 & -2 & 2 \\
$x_8$ & 0 & 1 & -1 & 0 & -1 & -2 & 0 & 1 & 2 & 7 \\
$x_1$ & 1 & -1 & -1 & 0 & 0 & -1 & 0 & 0 & 1 & 1 \\
\hline
\end{tabular}
\end{center}

$x_2$入基，比值检验：$\min\{2/1, 7/1\} = 2$，$x_7$出基。

第2次迭代后（基$B = \{x_2, x_8, x_1\}$）：

\begin{center}
\begin{tabular}{|c|rrrrrrrrr|r|}
\hline
B & $x_1$ & $x_2$ & $x_3$ & $x_4$ & $x_5$ & $x_6$ & $x_7$ & $x_8$ & $x_9$ & $b$ \\
\hline
 & 0 & 0 & 1 & 1 & 1 & 4 & 2 & 0 & -3 & -5 \\
\hline
$x_2$ & 0 & 1 & 0 & 1 & 0 & 2 & 1 & 0 & -2 & 2 \\
$x_8$ & 0 & 0 & -1 & -1 & -1 & -4 & -1 & 1 & 4 & 5 \\
$x_1$ & 1 & 0 & -1 & 1 & 0 & 1 & 1 & 0 & -1 & 3 \\
\hline
\end{tabular}
\end{center}

$x_9$入基，比值检验：$\min\{5/4\} = 5/4$，$x_8$出基。

第3次迭代后（基$B = \{x_2, x_9, x_1\}$），第一阶段达到最优：

\begin{center}
\begin{tabular}{|c|rrrrrrrrr|r|}
\hline
B & $x_1$ & $x_2$ & $x_3$ & $x_4$ & $x_5$ & $x_6$ & $x_7$ & $x_8$ & $x_9$ & $b$ \\
\hline
 & 0 & 0 & $\frac{1}{4}$ & $\frac{1}{4}$ & $\frac{1}{4}$ & 1 & $\frac{5}{4}$ & $\frac{3}{4}$ & 0 & $-\frac{5}{4}$ \\
\hline
$x_2$ & 0 & 1 & $-\frac{1}{2}$ & $\frac{1}{2}$ & $-\frac{1}{2}$ & 0 & $\frac{1}{2}$ & $\frac{1}{2}$ & 0 & $\frac{9}{2}$ \\
$x_9$ & 0 & 0 & $-\frac{1}{4}$ & $-\frac{1}{4}$ & $-\frac{1}{4}$ & -1 & $-\frac{1}{4}$ & $\frac{1}{4}$ & 1 & $\frac{5}{4}$ \\
$x_1$ & 1 & 0 & $-\frac{5}{4}$ & $\frac{3}{4}$ & $-\frac{1}{4}$ & 0 & $\frac{3}{4}$ & $\frac{1}{4}$ & 0 & $\frac{17}{4}$ \\
\hline
\end{tabular}
\end{center}

第一阶段$w^* = \dfrac{5}{4} > 0$，原问题\textbf{不可行}。

\textbf{结论：}原问题无可行解。

\end{document}
